% font: footnotesize
% TAB-26: generated by harness/analysis/paper_outputs.py; do not edit (every cell's source is in TAB-26.cells.json)
\begin{tabularx}{\linewidth}{>{\hsize=1.915\hsize\raggedright\hyphenpenalty=10000\exhyphenpenalty=10000\arraybackslash}X>{\hsize=2.233\hsize\raggedright\hyphenpenalty=10000\exhyphenpenalty=10000\arraybackslash}X>{\hsize=0.749\hsize\raggedleft\hyphenpenalty=10000\exhyphenpenalty=10000\arraybackslash}X>{\hsize=0.503\hsize\raggedleft\hyphenpenalty=10000\exhyphenpenalty=10000\arraybackslash}X>{\hsize=0.414\hsize\raggedleft\hyphenpenalty=10000\exhyphenpenalty=10000\arraybackslash}X>{\hsize=0.386\hsize\raggedright\hyphenpenalty=10000\exhyphenpenalty=10000\arraybackslash}X>{\hsize=0.799\hsize\raggedleft\hyphenpenalty=10000\exhyphenpenalty=10000\arraybackslash}X}
\toprule
What it is & Stamped key (of the collapsed interval; its field name follows where it is not ci95) & Stamped bootstrap 95\% accuracy interval (\%) & Stamped bootstrap 95\% interval of a paired difference (points) & Samples per puzzle (a~vote: its~$k$) & Puzzles right (of~the cell's puzzles; post hoc) & Exact 95\% accuracy interval, Clopper--Pearson (\%, post hoc) \\
\midrule
Length sweep (Deviation~2 rule), Opus-medium whole-answer, 6 steps, pooled & models.om.by\_n.6.whole\_answer.pooled & \mbox{[100.00, 100.00]} & --- & 1 & 20 of 20 & \mbox{[83.16, 100.00]} \\
Length sweep (Deviation~2 rule), Opus-medium per-step, 8 steps, pieces & models.om.by\_n.8.per\_step.nx1 & \mbox{[100.00, 100.00]} & --- & 1 & 20 of 20 & \mbox{[83.16, 100.00]} \\
Length sweep (Deviation~2 rule), Opus-medium per-step, 12 steps, chains & models.om.by\_n.12.per\_step.1xn & \mbox{[100.00, 100.00]} & --- & 1 & 19 of 19 & \mbox{[82.35, 100.00]} \\
Length sweep (Deviation~2 rule), Opus-medium per-step, 12 steps, pieces & models.om.by\_n.12.per\_step.nx1 & \mbox{[100.00, 100.00]} & --- & 1 & 20 of 20 & \mbox{[83.16, 100.00]} \\
Length sweep (Deviation~2 rule), Opus-medium per-step, 12 steps, pooled & models.om.by\_n.12.per\_step.pooled & \mbox{[100.00, 100.00]} & --- & 1 & 20 of 20 & \mbox{[83.16, 100.00]} \\
Length sweep (Deviation~2 rule), Opus-medium per-step, 16 steps, chains & models.om.by\_n.16.per\_step.1xn & \mbox{[100.00, 100.00]} & --- & 1 & 20 of 20 & \mbox{[83.16, 100.00]} \\
Length sweep (Deviation~2 rule), Opus-medium whole-answer, 16 steps, pooled & models.om.by\_n.16.whole\_answer.pooled & \mbox{[100.00, 100.00]} & --- & 1 & 20 of 20 & \mbox{[83.16, 100.00]} \\
Length sweep (Deviation~2 rule): a Pilot~1 marker & pilot1.om.l14\_only & \mbox{[100.00, 100.00]} & --- & not stamped & 5 of 5 & \mbox{[47.82, 100.00]} \\
\addlinespace
Free analysis 2: a Haiku-max vote-curve point & result.N=6.curve.8.1x6 & \mbox{[100.00, 100.00]} & --- & 9 & 40 of 40 & \mbox{[91.19, 100.00]} \\
Free analysis 2: a Haiku-max vote-curve point & result.N=6.curve.9.1x6 & \mbox{[100.00, 100.00]} & --- & 10 & 40 of 40 & \mbox{[91.19, 100.00]} \\
\addlinespace
Item~9e: the cascade's Opus-medium accuracy on its half of the puzzles & reporting.group\_b.item\_9e.shapes.6x1.5.even (field opus\_accuracy\_ci95) & \mbox{[100.00, 100.00]} & --- & 1 & 40 of 40 & \mbox{[91.19, 100.00]} \\
Item~9e: the cascade's Opus-medium accuracy on its half of the puzzles & reporting.group\_b.item\_9e.shapes.6x1.10.even (field opus\_accuracy\_ci95) & \mbox{[100.00, 100.00]} & --- & 1 & 40 of 40 & \mbox{[91.19, 100.00]} \\
Item~9e: the cascade's Opus-medium accuracy on its half of the puzzles & reporting.group\_b.item\_9e.shapes.6x1.20.even (field opus\_accuracy\_ci95) & \mbox{[100.00, 100.00]} & --- & 1 & 40 of 40 & \mbox{[91.19, 100.00]} \\
\bottomrule
\end{tabularx}
% note: Each stamped bootstrap interval has zero width because every puzzle in the cell was right, or every one wrong; the exact interval replaces it. Nothing here is a test, and nothing here changes a pre-registered reading.
% note: Exact interval: the two-sided 95\% Clopper--Pearson interval for the puzzles right out of the cell's puzzles, with the puzzle as the unit (the conservative unit: a puzzle with several samples, or a vote, counts once). It counts puzzles, not samples.
% note: Scope (free analysis 4, revision 3): the rows of this table and of its item~20 panel are every zero-width interval (a stamped interval whose two ends are equal) that free analysis 4 covers, under any interval field name, in the result files its scan reads (Study~1, main run; Study~1, keep-list analyses; Study~1b; Study~1b's samples, free analysis 2 (post hoc); Top-up run; Length sweep, original marking rule; Length sweep, scored under Deviation~2), except the 9 that are not reported (below).
% note: Cells: the previous revision's 105 cells, unchanged (the first version's 13, and 92 added by rule, all under the key ci95), and the 12 that this revision added by scanning every interval field by name: 3 are rows here (item~9e's Opus-medium accuracy on the even-indexed puzzles, stamped under opus\_accuracy\_ci95) and 9 are not reported.
% note: Not reported: this revision also computed 9 exact 90\% intervals (Clopper--Pearson, at the stamped interval's own level) for the fairness check's misapplied-rule rates, stamped under ci90; they are not reported (beyond the approved scope), and no table or figure prints them.
% note: Left out (each keeps its zero-width stamped interval and has no exact interval), with free analysis 4's counts and reasons: 1,061 zero-width intervals in Study~1, main run (``item~20 at $k$ below 40: ⟦FIG-02⟧ draws the vote at $k$~=~40 only, so these cells are printed nowhere''); 1 zero-width interval in Length sweep, original marking rule (``$q$ pools every piece-step over all lengths and sets (clusters stratified by set), so no single count of base puzzles carries it, and its $q^{N}$ bands are powers of its interval; no exact interval''); 6 zero-width intervals in Length sweep, original marking rule (``$q$'s $q^{N}$ bands: powers of $q$'s interval; $q$ has no exact interval''); 1 zero-width interval in Length sweep, scored under Deviation~2 (``as for the original rule's $q$''); 6 zero-width intervals in Length sweep, scored under Deviation~2 (``$q$'s $q^{N}$ bands: powers of $q$'s interval; $q$ has no exact interval'').
% note: Recount: the scan finds 1,191 zero-width intervals in these files, under any field name; 107 are behind the rows and the item~20 panel (a paired difference's one interval has a row for each of its two settings), 1,075 are left out and 9 are not reported.
% note: Scan: free analysis 4 reads every two-number field in these files; it counts 4,819 of them as intervals, those whose field name is one of: ci, cp or wilson followed by two digits (for example ci95), with or without a prefix that ends in an underscore; exactly `band'; or any name ending in \_bars, and the only one it does not count as an interval is quartiles (3 fields: the first and third quartiles of a spread, not an interval for an estimate).
% note: C10 rows: ``the stamped interval is the paired difference's \mbox{[0, 0]}; each arm's exact accuracy interval is given instead'' (free analysis 4's note, where its ``arm'' is a setting). The difference is Opus-low minus Opus-medium on pieces, in points, and is shown in its own column; a difference has no exact counterpart.
% note: Length sweep: $q$ is a model's per-step accuracy, the share of all its pieces right pooled over the sweep's lengths (4, 6, 8, 12 and 16 steps), and $q$ to the power $N$ is the reference for a chain of $N$ steps. The length sweep's by-length rows here are all Opus-medium's, and the inputs of $q$ are its pieces per-step rows (keys ending per\_step.nx1), here at 4, 6, 8, 12 and 16 steps under both marking rules, each over that length's puzzles; the mark ($q$'s input) in a row's name is the first version's own wording, so only its rows carry it (the rows added by rule are named by their family, model, measure, length and shape, and their keys show which are inputs of $q$); Opus-medium's $q$ itself, pooled over every length, has a zero-width stamped interval under each rule and no exact interval (the scope note gives free analysis 4's reason), while the other models' $q$ intervals have width.
% note: Labels: the counts and the exact intervals are post hoc presentation; each stamped interval keeps its source's label, given in its own column: the source's stamped label (for Study~1b's descriptive blocks, the block's name; for the length sweep, its stamped plan line; for the Pilot~1 markers the length-sweep files carry, exploratory, the label the plan gives Pilot~1's 24-step markers (they ran before Lock~2, with a different rule mix, sets and timing); for free analysis 2, its stamped label), or, where the source stamps none, the plan's label for the output that reports it (⟦FIG-02⟧ for item~20; ⟦TAB-11⟧ for Study~1b's per-test cells). Samples per puzzle: a vote row gives its $k$, the samples it votes over; a row with several single samples averages over them; either way the puzzle counts once in the count of puzzles right; a Pilot~1 marker's samples are not stamped in the length-sweep files.
% note: Item~24: its two rows are item~10's own-$k$ votes (checked equal in $k$, accuracy and cost), and the main analysis stamps no label for item~24, so they carry item~10's stamped label, ``exploratory (no Holm, no verdict)'', as ⟦TAB-05⟧ prints them.
% note: Names: Item~7 is a Study~1 reporting item, per-piece accuracy: the share of a sample's pieces that are right, averaged over the puzzle's designed samples and then over puzzles; Item~9e is a Study~1 reporting item, the offline escalate-to-Opus cascade: its key ends in the number of Haiku-medium samples in the cascade's vote ($k_c$) and then even, the even-indexed puzzles, where the cascade is scored; Opus-medium's accuracy there is one estimate, the same for every $k_c$, so these rows repeat it; C10 is the keep-list analysis of the paired difference in single-sample accuracy between two effort levels of one model; a vote of $k$ samples scores one answer per puzzle from those samples; a vote-curve point is the vote at one $k$; `pooled' in a row name or key means both shapes together, each puzzle counted once, not pooled over lengths; `chains' is the chain shape.
% note: Keys: a shape is written with the letter x, so 6x1 is the shape 6\ensuremath{\times}1; 1xn and nx1 are the chain and the pieces at the row's length; the number after `by\_n.' or after `N=' in a key is the row's length in steps; per\_step is per-step accuracy and whole\_answer whole-answer accuracy; a model is written as a tag (om: Opus-medium; ho: Haiku, thinking off), and tag@$N$ is that model at $N$ steps; points.j and curve.j are a point of a vote curve by its place in the list (its $k$ is in the samples column); Study~1b's descriptive.effort\_figure and descriptive.length\_curve keys are the copies its figures draw of a setting, or, under length\_curve.ladder, of a length-sweep block, each checked equal to what it copies; `own $k$ vs opus-med' (item~24) is the vote at the row's own cost-matched $k$ against Opus-medium; `own $k$ vs sonnet-med' (item~24) is the vote at the row's own cost-matched $k$ against Sonnet-medium; per\_form.B3.6x1.o is Test~B3's cell on the pieces for its Opus-medium setting (the test's per-shape breakdown); pilot1.om.l14\_only is a Pilot~1 result that the length-sweep files carry beside the sweep, stamped ``pilot-1, presets l14, before lock-2 (different rule mix, sets and timing)''.
% note: Words: each row's description is free analysis 4's own, in this paper's words where it used a word or a form of a name this paper no longer uses (setting for ``arm'', sample for ``try'', puzzle for ``base puzzle'', cost-matched for ``equal-money'', ``Haiku, thinking off'' for ``Haiku thinking-off'', ``-sample vote'' after a vote's $k$ for ``vote@$k$''), each name starting with a capital, and a length-sweep row added by rule named in full from its key (its family, then the model, measure, length and shape); a stamped key that holds such a word is printed as stamped, in quotation marks.
% note: Values stamped as shares are printed in \%.
